% ── §18 Conclusion ──
\section{Conclusion}
\label{sec:conclusion}

We have demonstrated that the formal chain from biological sequences through
Gray-code encoding, Karnaugh-map construction, Quine--McCluskey minimisation,
and prime-implicant extraction can be made machine-provable in Lean~4, and
that the resulting Boolean circuits carry genuine structural information when
applied to protein multiple sequence alignments.  The central achievement is
not that these circuits match state-of-the-art contact predictors in absolute
precision---they do not, and we report this gap honestly.  Rather, it is that
the Karnaugh-map cell-rate prior carries orthogonal structural information
demonstrable by rank-fusion with continuous DCA scores, and that every step of
the encoding pipeline is verified by machine-checked proofs rather than by
empirical observation alone.

The framework's principal value lies in three areas.  First, it provides a
formally verified foundation upon which future coevolution analyses can build,
eliminating an entire class of encoding errors.  Second, it produces
interpretable sum-of-products rules that biologists can inspect, validate
against known structural principles (such as hydrophobic-core packing and
salt-bridge formation), and potentially use as design constraints in protein
engineering.  Third, its cross-protein transfer capability demonstrates that
the learned rules capture general principles of protein architecture rather
than family-specific idiosyncrasies.

The path from here to full structure reconstruction remains long.  It requires
closing the precision gap between Boolean circuit features and continuous
coupling scores, incorporating phylogenetic correction into the discrete
framework without sacrificing interpretability, and scaling exact minimisation
to larger feature spaces.  We believe the formal verification infrastructure
established here provides the correct foundation for that journey.
